\section{Static instance records}\label{si:static}

\subsection{One-dimensional instances}
All on $U=[0,1]$; values exact (checks S1.1--S1.6).
\begin{center}\footnotesize
\begin{tabularx}{\textwidth}{llcc>{\raggedright\arraybackslash}X}
\toprule
instance & rows & $V^-$ & maximiser & dual / certificate\\
\midrule
two-floor & $\tfr25-u,\ u-\tfr35$ & $-\tfr1{10}$ & $\tfr12$ & $\mu^\ast=(\tfr12,\tfr12)$, pooled row $\equiv-\tfr1{10}$, margin $\tfr1{10}$\\
feasible contrast & $\tfr25-u,\ u-\tfr15$ & $+\tfr1{10}$ & $\tfr3{10}$ & no certificate (dual value $+\tfr1{10}$)\\
Isaacs & $1-2u,\ 2u-1$ & $0$ & $\tfr12$ & dual value $0$; upper value $V^+=1$\\
parity ($U=\{0,1\}$) & $2u-1,\ 1-2u$ & $-1$ over $U$ & -- & value $0$ over $\conv U$; no certificate\\
\bottomrule
\end{tabularx}
\end{center}
The two-floor instance is the active bundle of two states with floors at $\tfr25$ and $\tfr35$ and drifts of opposite sign: state $1$ needs $u\le\tfr25$, state $2$ needs $u\ge\tfr35$. The Isaacs instance separates the lower value (common action $u=\tfr12$, both rows at $0$) from the upper value (each row alone is satisfied with margin $1$ at $u=0$ or $u=1$).

\subsection{The tight family (main Theorem~\mref{thm:helly}(ii))}
Rows $u_i-1\ge0$ $(i\le k)$ and $(k-1)-\sum_iu_i\ge0$ on $U=[0,k]^k$ (checks S2.1--S2.5).
\begin{center}\footnotesize
\begin{tabular}{cccccl}
\toprule
$k$ & rows & every $k$-subfamily feasible & lower value & uniform margin & plant\\
\midrule
1 & 2 & yes & $-\tfr12$ & $\tfr12$ & $\dot x_1=-1+u_1,\ \dot x_2=-u_1$\\
2 & 3 & yes & $-\tfr13$ & $\tfr13$ & $\dot x_i=-1+u_i,\ \dot x_3=1-u_1-u_2$\\
3 & 4 & yes & $-\tfr14$ & $\tfr14$ & $\dot x_i=-1+u_i,\ \dot x_4=2-\sum u_i$\\
4 & 5 & yes & $-\tfr15$ & $\tfr15$ & $\dot x_i=-1+u_i,\ \dot x_5=3-\sum u_i$\\
\bottomrule
\end{tabular}
\end{center}
The lower value is attained at $u=(\tfr{k}{k+1},\dots,\tfr{k}{k+1})$; the plant realisation evaluates $\nabla q_j^\top f$ at $x=0$ with $q_j(x)=x_j$, all $k+1$ constraints active, and the script checks that the rows read off the plant coincide with the listed rows.

\subsection{The benchmark atoms}
At the corner $x^{(1)}=(2,Y-2)$ the active constraint is $q_1$ and the row is $\dot x_1=r_1\cdot2\cdot(1-\tfr2{10})-\tfr{2(Y-2)}{20}-u_1+d_1$; at $x^{(2)}=(Y-2,2)$ the row is $\dot x_2=r_2\cdot2\cdot\tfr45-\tfr{2(Y-2)}{20}-u_2+d_2$. With both shock values kept the bundle has four atoms,
\[
\bigl(x^{(1)},1,\mp\tfr1{10}\bigr),\qquad \bigl(x^{(2)},2,\mp\tfr1{10}\bigr),
\]
with constants $\capi1(Y)$ and $\capi1(Y)+\tfr15$ for the first pair, $\capi2(Y)$ and $\capi2(Y)+\tfr15$ for the second, and coefficient vectors $(-1,0)$ and $(0,-1)$ respectively. The exact simplex programme over the four atoms returns the same value as over the two worst-shock atoms and charges only the latter: $\lambda=(\tfr12,0,\tfr12,0)$ (check B2.3). Caps and margins at the quoted aggregates (checks B1.1--B2.2):
\begin{center}\footnotesize
\begin{tabularx}{\textwidth}{cccccc>{\raggedright\arraybackslash}X}
\toprule
$Y$ & $\capi1$ & $\capi2$ & $\capi1+\capi2$ & margin $\varepsilon^\ast(Y)$ & maximiser & status\\
\midrule
$4$ & $\tfrac{13}{10}$ & $\tfrac{49}{50}$ & $\tfrac{57}{25}$ & $\tfrac{7}{50}$ & $(\tfrac{29}{25},\tfrac{21}{25})$ & viable\\
$5$ & $\tfrac{6}{5}$ & $\tfrac{22}{25}$ & $\tfrac{52}{25}$ & $\tfrac{1}{25}$ & $(\tfrac{29}{25},\tfrac{21}{25})$ & viable; $(\tfrac65,\tfrac45)$ also safe\\
$\tfrac{27}{5}$ & $\tfrac{29}{25}$ & $\tfrac{21}{25}$ & $2$ & $0$ & $(\tfrac{29}{25},\tfrac{21}{25})$ & boundary\\
$6$ & $\tfrac{11}{10}$ & $\tfrac{39}{50}$ & $\tfrac{47}{25}$ & $-\tfrac{3}{50}$ & -- & obstructed, certificate $(\tfrac12,\tfrac12)$\\
$\tfrac{69}{5}$ & $\tfrac{8}{25}$ & $0$ & $\tfrac{8}{25}$ & $-\tfrac{21}{25}$ & -- & corner 2 at its cap\\

\bottomrule
\end{tabularx}
\end{center}
The maximiser $(\tfr{29}{25},\tfr{21}{25})$ is independent of $Y$ because the caps shift together; the margin sweep on $Y\in[4,7]$ in steps of $\tfr1{20}$ agrees with $(27-5Y)/50$ at every point and the dual measure is $(\tfr12,\tfr12)$ at every point (check B2.1).

\subsection{The demand-floor law}
$Y^\ast(D)=\min(\tfr{77}5-5D,\tfr{69}5)$ for $D\le\tfr{57}{25}$ (checks B4.1--B4.4). Selected values:
\begin{center}\footnotesize
\begin{tabular}{cccl}
\toprule
$D$ & $\tfr{77}5-5D$ & $Y^\ast(D)$ & binding obstruction at $Y^\ast$\\
\midrule
$0$ & $\tfrac{77}{5}$ & $\tfrac{69}{5}$ & physical: $\operatorname{cap}_2=0$ at corner 2\\
$\tfrac{8}{25}$ & $\tfrac{69}{5}$ & $\tfrac{69}{5}$ & both regimes meet\\
$\tfrac{1}{2}$ & $\tfrac{129}{10}$ & $\tfrac{129}{10}$ & epistemic + institutional\\
$1$ & $\tfrac{52}{5}$ & $\tfrac{52}{5}$ & epistemic + institutional\\
$2$ & $\tfrac{27}{5}$ & $\tfrac{27}{5}$ & epistemic + institutional (main text)\\
$\tfrac{57}{25}$ & $4$ & $4$ & last viable fibre is the corner fibre $Y=4$\\
$>\tfrac{57}{25}$ & $<4$ & none & every fibre obstructed\\

\bottomrule
\end{tabular}
\end{center}
The slope $-5$ is $-(2\cdot\text{comp}\cdot\text{floor})^{-1}$ with comp $=\tfr1{20}$ and floor $=2$: the cap sum decreases by $\tfr15$ per unit of $Y$, so one unit of demand floor is worth five units of aggregate stock. The script checks, for $D$ on the grid $\tfr8{25},\tfr9{25},\dots,\tfr{57}{25}$, that the margin is zero at $Y^\ast(D)$, positive at $Y^\ast(D)-\tfr1{100}$ and negative at $Y^\ast(D)+\tfr1{100}$, and that for $D=0$ the sign change happens at $\tfr{69}5$.
